\documentclass[11pt]{article}
\usepackage[T1]{fontenc}
\usepackage{lmodern,amsmath,amssymb,amsthm,booktabs,mathtools}
\usepackage[margin=0.86in]{geometry}
\usepackage[colorlinks=true,allcolors=blue!60!black]{hyperref}
\usepackage{xcolor}
\newtheorem{theorem}{Theorem}
\newtheorem{lemma}[theorem]{Lemma}
\newtheorem{proposition}[theorem]{Proposition}
\newcommand{\Z}{\mathbb Z}
\title{All prime exponents of a three-prime divisor-moment fixed point must be at least four}
\author{Blake Stone\\\small Developed with OpenAI Codex}
\date{November 26, 2026}
\begin{document}
\maketitle
\begin{abstract}
Let $S(n)$ be the sum of $d\tau(d)$ over the proper positive divisors of $n$.
We prove that if $S(n)=n$ has a solution with three distinct prime bases, all exponents are at least four. For ordered exponents $a\ge b\ge c$, we also prove $c\le25$, $b\le10c+7$, and $a\le(3c+4)b+3c$. In the case $c=1$ the sharper bounds $b\le16$ and $a\le7b+3$ permit an exact enumeration. For $c=2$ a modular sieve leaves only two exponent triples. We also give a uniform obstruction on
individual exponents at every odd prime.
\end{abstract}

\section{The equation and local obstructions}
Jaroslav Krizek proposed that $21$ is the unique positive fixed point of
$S$ in OEIS A245211 \cite{source}. Earlier results include the
prime-factor lower bound and local parity restrictions \cite{earlier,sugui}.
The preceding note \cite{companion} classifies inputs with at most two
prime factors and excludes $p^a q r$. We extend its rational
approximation method to two variable exponents.

Write $n=\prod p_i^{e_i}$, $D=\tau(n)=\prod(e_i+1)$, and set
\[
 H_e(t)=\sum_{j=0}^e(j+1)t^j,\qquad F_e(t)=t^{-e}H_e(t),\qquad
 T_e(t)=F_e(t)-(e+1).
\]
Multiplicativity of the full divisor sum and the contribution of the
proper divisor $n/p_i$, together with the separate term $1$, give
\begin{equation}
 \prod_i F_{e_i}(p_i)=D+1,\qquad
 p_i>e_iD/(e_i+1)\ge D/2.\label{eq:main}
\end{equation}
Here and below the prime lower bound is used for inputs with at least two
distinct prime factors. For $e\ge1$ and $t>1$, summing the finite
geometric progression gives
\begin{equation}
 \frac{e}{t}\le T_e(t)=\frac{e}{t-1}-\frac{1-t^{-e}}{(t-1)^2}
 <\frac{e}{t-1},\qquad
 F_e(t)=\frac{(e+1)t^2-(e+2)t+t^{-e}}{(t-1)^2}.\label{eq:tail}
\end{equation}

\begin{proposition}\label{prop:local}
For an odd prime $\ell$, $H_e$ vanishes on all units modulo $\ell$
exactly when $\ell(\ell-1)$ divides $e+1$ or $e+2$. No exponent of a
fixed point can satisfy the first divisibility. In the second case a
fixed point requires $\prod_{j\ne i}(e_j+1)\equiv1\pmod\ell$ when
$e=e_i$.
\end{proposition}
\begin{proof}
At $t=1$ the value is $(e+1)(e+2)/2$. If $s$ is the member of
$\{e+1,e+2\}$ divisible by $\ell$, then
$(t-1)^2 H_e(t)\equiv1-t^s\pmod\ell$ by \eqref{eq:tail}.
A generator of the cyclic group of units gives the condition
$\ell-1\mid s$; the same condition suffices on all units.
In either divisibility case $e\ge\ell(\ell-1)-2\ge\ell$.
The prime base of this exponent exceeds $e$ by \eqref{eq:main}.
A prime power alone cannot be a fixed point, as its proper moment
is $1$ modulo its base.
Thus the moment of the designated prime power is divisible by $\ell$.
If some other prime base were $\ell$, its lower bound
$p_j>e_j\prod_{h\ne j}(e_h+1)\ge e+1$ rules this out.
Clearing the bases in \eqref{eq:main} now gives $D+1\equiv0\pmod\ell$.
This contradicts $\ell\mid e+1$ and gives the stated congruence when
$\ell\mid e+2$.
\end{proof}

For example, no fixed point has any exponent $5$ modulo $6$, and an
exponent $4$ modulo $6$ requires the product of the other
exponent increments to be $1$ modulo $3$.

\section{Bounds on three variable exponents}
For any initial list of chosen prime bases, define $R$ as $D+1$ minus
the product in \eqref{eq:main} with each remaining factor $F_e(t)$
replaced by $e+1$. If a factor remains then $R>0$ and
\begin{equation}
 \frac{R}{D+1}=1-\prod_{\mathrm{remaining}}\frac{e+1}{F_e(t)}
 <\sum_{\mathrm{remaining}}\frac1{t-1}.\label{eq:res}
\end{equation}
Indeed $1-\prod(1-u_i)\le\sum u_i$ for $u_i\in[0,1]$, and
$1-(e+1)/F_e(t)<1/(t-1)$ by \eqref{eq:tail}.

\begin{lemma}\label{lem:single}
Suppose $p^a q^b z^c$ is fixed with distinct prime bases and $q<z$.
Then $q<p^3$ and $a<3bc+3b+4c+1$.
\end{lemma}
\begin{proof}
Let $K=(b+1)(c+1)$, so $D=(a+1)K$ and $p-1\ge aK$.
With $p$ chosen, \eqref{eq:tail} gives
\[
 R_1=1-KT_a(p)\ge\frac{K(1-p^{-a})}{(p-1)^2}
 \ge\frac{K}{p(p-1)}.
\]
Hence $q-1<2(D+1)p(p-1)/K$ by \eqref{eq:res}. Since
$2(D+1)/K=2a+2+2/K\le aK+1\le p$, this gives $q<p^3$.

Suppose that $a\ge3bc+3b+4c+1$, so $a\ge4$ and $a\ge3b+1$.
For the two chosen bases $p,q$, the positive residual is
$R_2=D+1-(c+1)F_a(p)F_b(q)$. Clearing the denominator of the main
term for $F_a(p)$ in \eqref{eq:tail} gives an integer $J$ such that
\[
 (p-1)^2 q^bR_2=J-\frac{(c+1)H_b(q)}{p^a}.
\]
Its positivity forces $J\ge1$. At any prime base of a fixed point
$F_e(t)<e+2$ by \eqref{eq:main} and \eqref{eq:tail}.
Since $p>aK\ge4K$, it follows that
\[
 2(c+1)H_b(q)<4Kq^b<p^{3b+1}\le p^a.
\]
Thus $R_2>1/[2(p-1)^2q^b]$. Equation \eqref{eq:res} gives
$z<1+2(D+1)(p-1)^2q^b<p^{3b+4}$, using $D+1<2p$ and $p>4$.
The original integer identity and $H_a(p)\equiv1\pmod p$ now imply
\[
 p^a\le H_b(q)H_c(z)<(b+2)(c+2)q^b z^c
 <p^{1+3b+(3b+4)c}.
\]
Here $(b+2)(c+2)\le(9/4)K<p$. This contradicts our assumption on $a$.
\end{proof}

\begin{lemma}\label{lem:both}
If $p^a q^b r$ with distinct prime bases is fixed and $a\ge b$, then
$b\le16$ and $a\le7b+3$.
\end{lemma}
\begin{proof}
The upper bound on $a$ follows from Lemma~\ref{lem:single} in both
orders of the other bases. For the other assertion assume $b\ge17$.
List the three bases as $x<y<z$, in an arbitrary order of the exponents.
Then $x>D/2=(a+1)(b+1)\ge324$ and $y<x^3$.
For a base $t$ of exponent $e\ge b$ replace $F_e(t)$ by
\[
 F'_e(t)=F_e(t)-\frac1{t^e(t-1)^2}.
\]
Keep the simple factor $F_1(t)$ exact. For the first $s$ bases chosen in
order, write $R_s$ for the preceding residual and $J_s$ for the residual
with these replacements. Every $J_s$ is positive: for $s<3$ this
follows from $R_s>0$, and at $s=3$ at least one factor strictly decreases.
Subtracting the two products term by term, the remaining factors have
product at most $D+1$. Hence
\begin{equation}
 0\le J_s-R_s \le \frac{2(D+1)}{x^b(x-1)^2}
 <\frac5{x^{b+1}}.\label{eq:error}
\end{equation}
Assign $d_t=1$ to the simple base and $d_t=2$ to both other bases.
The denominator of $J_s$ is at most $Q_s=\prod_{t\ \mathrm{chosen}}t^{d_t}$,
so $J_s\ge1/Q_s$. As $Q_2\le x^2 y^2<x^8$, \eqref{eq:error} gives
$R_2>1/(2Q_2)$. Equation \eqref{eq:res} and integrality then give
$z\le4xQ_2$.

We claim $Q_3\le16 x^{17}$. If $z$ is simple then
$Q_3=Q_2 z\le4xQ_2^2<4x^{17}$. If $y$ is simple then $Q_2=x^2y<x^5$,
so $Q_3\le16x^2Q_2^3<16x^{17}$.
If $x$ is simple, $2x+1>D+1$ is a divisor of $(D+1)n$. Hence some
prime divisor of $n$ divides $2x+1$. It exceeds $x$, so must equal
$2x+1$ itself. In particular $y\le2x+1<3x$.
Now $Q_2=xy^2<9x^3$, whence $Q_3<16\cdot9^3x^{11}<16x^{17}$.
But $R_3=0$ and $5/x^{b+1}<1/(16x^{17})$ since $x>80$.
This contradicts $J_3\ge1/Q_3$ and \eqref{eq:error}.
\end{proof}

\section{Finite determination}
\begin{theorem}
There is no fixed point of $S$ with exactly three distinct prime
divisors if any one of them has exponent one.
\end{theorem}
\begin{proof}
By Lemma~\ref{lem:both} it suffices to test $1\le b\le16$ and
$b\le a<7b+4$. The following exact procedure exhausts the possibilities.
At any modulus $m<D/2$ all prime bases are units. Enumerate the
values of $F_e(u)$ on the units for each exponent and discard a
pattern if their products omit $D+1$. The code lists the small
moduli used. For each ordered assignment $(i,j,k)$ of the exponents
to $x<y<z$, test primes $iD/(i+1)<x\le1+3(D+1)$ as allowed by
\eqref{eq:main} and \eqref{eq:res}. Repeat the modular test with
$x$ fixed. With $A=H_i(x)$, let $g=(D+1)x^i-(j+1)(k+1)A$.
We need $g>0$ and
\[
 x<y\le Y=1+\left\lfloor2(D+1)x^i/g\right\rfloor .
\]
The second inequality follows from \eqref{eq:res}. As $\gcd(A,x)=1$,
the integer $B=A/\gcd(A,D+1)$ divides $y^jz^k$. Also $B>1$ since
$A\ge2x+1>D+1$.

If a prime factor $t$ of $B$ is known, one of $y,z$ must be $t$.
Test both assignments. When $i=1$ the proof of Lemma~\ref{lem:both}
already forces one of $y,z$ to be $2x+1$. The argument there only
uses $x>D/2$.
Otherwise seek a prime factor of $B$ through $Y$, permitting a
larger preliminary trial bound. If none is found, then $y$ cannot
divide $B$, so $B=z^h$ for some $1\le h\le k$. Test every
integer root from these finitely many powers. Thus each surviving
branch supplies one of the two remaining bases.
For the other, call the fixed base $u$ with exponent $f$, the
unknown $v$ with exponent $e$, and put
$W=H_i(x)H_f(u)$ and $G=(D+1)x^iu^f-(e+1)W$.
At a solution $G/W=T_e(v)>0$. The two bounds in \eqref{eq:tail}
give exactly one candidate $v=\lceil eW/G\rceil$. Test the integer
moment equation, retaining $x<y<z$ with odd bases.

\begin{center}
\begin{tabular}{lrrrrrrr}
\toprule
Smaller exponent $b$ & 1&4&8&9&12&13&16\\
Admissible $a$ after the modular test &4&2&10&19&13&20&3\\
\bottomrule
\end{tabular}
\end{center}
Every other $b$ has zero admissible exponents. There are $15{,}146$
ordered smallest-prime rows after sieving. The largest of the 42
extended prime-divisor intervals ends at $254{,}942{,}533$.
Among $18{,}675$ single-candidate completions, the only ordered
odd-base equality retained is $(x,y,z)=(7,15,31)$ with all
exponents one. Its middle base is composite. This establishes the
theorem.
\end{proof}

The accompanying Python and C++ programs carry out the finite proof.
The independent checks use a second segmented prime sieve, literal
moment divisibility and integer bisection for the missing base.
\section{Exponents two and three, and a global box}
\begin{theorem}
Suppose a fixed point has exactly three prime divisors, with exponents
$a\ge b\ge c$. Then
\[
 4\le c\le25,\qquad b\le10c+7,\qquad a\le(3c+4)b+3c.
\]
In particular every prime dividing such a number divides it to at
least the fourth power.
\end{theorem}
\begin{proof}
For the upper bound on $a$, apply Lemma~\ref{lem:single} with the two
smaller exponents in either order. Continue to write $x<y<z$ for the
ordered bases, so $y<x^3$. First suppose $c\ge26$.
Truncate all three factors using \eqref{eq:tail}, as in
Lemma~\ref{lem:both}. Assign $d_t=2$ to every base, and keep the
notation $J_s,R_s,Q_s$ from that proof. Positivity and the denominators
give $J_s\ge1/Q_s$. Since $D+1<2x$ and $x>D/2>200$, the total error
in any prefix is less than $7/x^{c+1}$. We get $Q_2<x^8$,
$R_2>1/(2Q_2)$, and $z\le4xQ_2$ as before. Hence
$Q_3\le16x^2Q_2^3<16x^{26}$. But $7/x^{c+1}<1/(16x^{26})$,
contradicting $R_3=0$. Thus $c\le25$.

For $c\ge2$ keep one factor of exponent $c$ exact, assigning
its base $d_t=c$; truncate the other two, assigning each $d_t=2$.
Suppose $b\ge10c+8$. Now $J_s-R_s<5/x^{b+1}$ and
$Q_2<x^{2+3c}$, so $R_2>1/(2Q_2)$ and $z\le4xQ_2$.
If the exact base is $z$, then $Q_2<x^8$ and
$Q_3\le4^c x^{8+9c}$. If it is $y$, then
$Q_2<x^{2+3c}$ and $Q_3\le16x^{8+9c}$. If it is $x$,
then $Q_2<x^{c+6}$ and $Q_3\le16x^{3c+20}$.
Since $c\ge2$ these all give $Q_3\le4^c x^{8+9c}<x^{8+10c}$.
The error bound again contradicts $R_3=0$. This proves the bound on $b$.

When $c=2$, the denominator exponent for its exact factor is two
as well. If $b\ge26$ then $Q_2<x^8$ and $Q_3<16x^{26}$ by the
first calculation above, with tail error less than $5/x^{b+1}$.
Thus $b\le25$ and $a\le10b+6$. The modular test in Section~3,
with the additional moduli $64$ and $81$, leaves precisely
$(6,2,2)$ and $(50,6,2)$. A separate direct enumeration tries each
permutation of each triple, every eligible prime $x$, and every eligible
prime $y$ through the bound $Y$ of Section~3. The largest actual
second-prime bound is $342{,}300$. Testing the missing last candidate
against the full moment equation gives no equality in $47{,}103$
integral tests. This excludes $c=2$. Finally an exponent three
makes its local moment even at any odd base, whereas $D+1$ is odd.
All three bases are odd by \eqref{eq:main}, giving a contradiction.
Section~3 excluded $c=1$.
\end{proof}
The conjecture is still open for the remaining exponent triples and
for unrestricted numbers of prime factors.

\small\begin{thebibliography}{9}
\bibitem{source} J. Krizek, \href{https://oeis.org/A245211}{OEIS A245211}
(2014).
\bibitem{earlier} The agentic-conjectures contributors,
\href{https://github.com/umaia1234/agentic-conjectures/blob/78a56edb19eadb6a23f5f791273f016e0b3e0f50/problems/oeis-a245211/DETAILS.md}{Weighted proper-divisor sum},
especially the prime-factor lower bound, parity and the last-prime test.
\bibitem{sugui} M. A. Sugui,
\href{https://doi.org/10.5281/zenodo.24437886}{Formal Verification of Parity Obstructions and Fixed-Point Exclusions for Proper-Divisor Moments} (2026).
\bibitem{companion} B. Stone with OpenAI Astra,
\href{https://github.com/blakestone-x/math/blob/a80ee14b88981065189113eace98e1b8fba58c28/A245211/A245211-research.pdf}{Fixed points of the weighted proper-divisor sum in OEIS A245211} (2026).
\end{thebibliography}
\end{document}
